\documentclass[10pt,a4paper]{amsart}
\usepackage{amsmath,amssymb,amsthm,graphicx}
\usepackage{verbatim}
\newenvironment{reference}{\comment}{\endcomment}
%\newenvironment{reference}{}{}
\newenvironment{slogan}{\comment}{\endcomment}
\newenvironment{history}{\comment}{\endcomment}

% For commutative diagrams we use Xy-pic
\usepackage[all]{xy}

% We use 2cell for 2-commutative diagrams.
\xyoption{2cell}
% Theorem environments.
%
\theoremstyle{plain}
\newtheorem{theorem}[subsection]{Theorem}
\newtheorem{proposition}[subsection]{Proposition}
\newtheorem{lemma}[subsection]{Lemma}

\theoremstyle{definition}
\newtheorem{definition}[subsection]{Definition}
\newtheorem{example}[subsection]{Example}
\newtheorem{exercise}[subsection]{Exercise}
\newtheorem{situation}[subsection]{Situation}

\theoremstyle{remark}
\newtheorem{remark}[subsection]{Remark}
\newtheorem{remarks}[subsection]{Remarks}

\numberwithin{equation}{subsection}

% Macros
%
\def\lim{\mathop{\mathrm{lim}}\nolimits}
\def\colim{\mathop{\mathrm{colim}}\nolimits}
\def\Spec{\mathop{\mathrm{Spec}}}
\def\Hom{\mathop{\mathrm{Hom}}\nolimits}
\def\Ext{\mathop{\mathrm{Ext}}\nolimits}
\def\SheafHom{\mathop{\mathcal{H}\!\mathit{om}}\nolimits}
\def\SheafExt{\mathop{\mathcal{E}\!\mathit{xt}}\nolimits}
\def\Sch{\mathit{Sch}}
\def\Mor{\mathop{\mathrm{Mor}}\nolimits}
\def\Ob{\mathop{\mathrm{Ob}}\nolimits}
\def\Sh{\mathop{\mathit{Sh}}\nolimits}
\def\NL{\mathop{N\!L}\nolimits}
\def\CH{\mathop{\mathrm{CH}}\nolimits}
\def\proetale{{pro\text{-}\acute{e}tale}}
\def\etale{{\acute{e}tale}}
\def\QCoh{\mathit{QCoh}}
\def\Ker{\mathop{\mathrm{Ker}}}
\def\Im{\mathop{\mathrm{Im}}}
\def\Coker{\mathop{\mathrm{Coker}}}
\def\Coim{\mathop{\mathrm{Coim}}}

% Boxtimes
%
\DeclareMathSymbol{\boxtimes}{\mathbin}{AMSa}{"02}

%
% Macros for moduli stacks/spaces
%
\def\QCohstack{\mathcal{QC}\!\mathit{oh}}
\def\Cohstack{\mathcal{C}\!\mathit{oh}}
\def\Spacesstack{\mathcal{S}\!\mathit{paces}}
\def\Quotfunctor{\mathrm{Quot}}
\def\Hilbfunctor{\mathrm{Hilb}}
\def\Curvesstack{\mathcal{C}\!\mathit{urves}}
\def\Polarizedstack{\mathcal{P}\!\mathit{olarized}}
\def\Complexesstack{\mathcal{C}\!\mathit{omplexes}}
% \Pic is the operator that assigns to X its picard group, usage \Pic(X)
% \Picardstack_{X/B} denotes the Picard stack of X over B
% \Picardfunctor_{X/B} denotes the Picard functor of X over B
\def\Pic{\mathop{\mathrm{Pic}}\nolimits}
\def\Picardstack{\mathcal{P}\!\mathit{ic}}
\def\Picardfunctor{\mathrm{Pic}}
\def\Deformationcategory{\mathcal{D}\!\mathit{ef}}
\usepackage{cleveref}
\crefname{originallemma}{original lemma}{original lemmas}
\Crefname{originallemma}{Original lemma}{Original lemmas}
\makeatletter
\newlabel{descent-lemma-properties-quasi-coherent}{{0GNB}{source}{}{sourcepointer-1}{}}
\newlabel{descent-lemma-properties-quasi-coherent@cref}{{[originallemma][0][]0GNB}{[1][0][]source}{}{}{}}
\makeatother
\crefname{originalequation}{original equation}{original equations}
\Crefname{originalequation}{Original equation}{Original equations}
\makeatletter
\newlabel{equation-localize-at-scheme-ringed}{{04M1}{source}{}{sourcepointer-2}{}}
\newlabel{equation-localize-at-scheme-ringed@cref}{{[originalequation][0][]04M1}{[1][0][]source}{}{}{}}
\makeatother
\crefname{originalequation}{original equation}{original equations}
\Crefname{originalequation}{Original equation}{Original equations}
\makeatletter
\newlabel{equation-restrict-modules}{{03LW}{source}{}{sourcepointer-3}{}}
\newlabel{equation-restrict-modules@cref}{{[originalequation][0][]03LW}{[1][0][]source}{}{}{}}
\makeatother
\crefname{originalreference}{original reference}{original references}
\Crefname{originalreference}{Original reference}{Original references}
\makeatletter
\newlabel{lemma-characterize-quasi-coherent}{{03M0}{source}{}{sourcepointer-4}{}}
\newlabel{lemma-characterize-quasi-coherent@cref}{{[originalreference][0][]03M0}{[1][0][]source}{}{}{}}
\makeatother
\crefname{originalreference}{original reference}{original references}
\Crefname{originalreference}{Original reference}{Original references}
\makeatletter
\newlabel{lemma-pullback-quasi-coherent}{{01BG}{source}{}{sourcepointer-5}{}}
\newlabel{lemma-pullback-quasi-coherent@cref}{{[originalreference][0][]01BG}{[1][0][]source}{}{}{}}
\makeatother
\crefname{originalsection}{original section}{original sections}
\Crefname{originalsection}{Original section}{Original sections}
\makeatletter
\newlabel{section-localize}{{00XZ}{source}{}{sourcepointer-6}{}}
\newlabel{section-localize@cref}{{[originalsection][0][]00XZ}{[1][0][]source}{}{}{}}
\makeatother
\crefname{originalsection}{original section}{original sections}
\Crefname{originalsection}{Original section}{Original sections}
\makeatletter
\newlabel{section-properties-modules}{{05VR}{source}{}{sourcepointer-7}{}}
\newlabel{section-properties-modules@cref}{{[originalsection][0][]05VR}{[1][0][]source}{}{}{}}
\makeatother
\crefname{originallemma}{original lemma}{original lemmas}
\Crefname{originallemma}{Original lemma}{Original lemmas}
\makeatletter
\newlabel{sites-modules-lemma-internal-hom-restriction}{{0E8H}{source}{}{sourcepointer-8}{}}
\newlabel{sites-modules-lemma-internal-hom-restriction@cref}{{[originallemma][0][]0E8H}{[1][0][]source}{}{}{}}
\makeatother
\crefname{originalsection}{original section}{original sections}
\Crefname{originalsection}{Original section}{Original sections}
\makeatletter
\newlabel{sites-modules-section-localize}{{03DH}{source}{}{sourcepointer-9}{}}
\newlabel{sites-modules-section-localize@cref}{{[originalsection][0][]03DH}{[1][0][]source}{}{}{}}
\makeatother
\title{The category of quasi-coherent sheaves on an algebraic space is closed under sums, colimits, kernels, tensor products and internal homs}
\author{The Stacks Project Authors}
\date{stacks-project:spaces-properties.tex:local-snapshot-20261012; GNU FDL 1.2 or later}
\begin{document}
\maketitle
\noindent\textit{Editorial note.} All delivered mathematical text is the author's own, copied from the pinned Stacks Project chapter file spaces-properties.tex; no statement or proof step was rewritten, reordered or removed.\par\medskip
\noindent\textit{Editorial note.} The counted proof is the author's complete argument: the reduction through a surjective etale cover by a scheme with the detection criterion for quasi-coherence and the commutation of restriction with limits and colimits, the four additive properties transported through the commutation with restriction, and the treatment of tensor products and internal homs through the commutation with localization.\par\medskip
\noindent\textit{Editorial note.} The label lemma-properties-quasi-coherent names the counted statement, and the proof environment immediately following it in the source is delivered verbatim as its proof.\par\medskip
\noindent\textit{Editorial note.} The counted statement is carried by the author's own lemma environment in the chapter source.\par\medskip
\noindent\textit{Editorial note.} Pre-open gate: python operations/agent-authoring/tools/precard\_gate.py --source spaces-properties.tex --label lemma-properties-quasi-coherent --start 5657 --end 5772 --local-id 28042 returned RESULT=OK (L28042 free, occupancy FREE, public interval clear). Candidate taken from sources/stacks-spaces-properties/candidates-halfA.json (fourth pool, pre-assigned by the Lead on 2026-10-13).\par\medskip
\noindent\textit{Editorial note.} The statement and proof are copied from the pinned chapter file; the cited inputs are kept as external pointers with their canonical Stacks tags.\par\medskip
\noindent\textit{Allowed standard prerequisites (editorial).}\begin{itemize}
\item The corresponding statement for schemes. Descent, the lemma listing the same six properties of the category of quasi-coherent sheaves on a scheme. Supplies the six properties in the case where the algebraic space is a scheme.
\item Detection of quasi-coherence after restriction to an etale cover. Same chapter, the lemma that a module on an algebraic space is quasi-coherent exactly when its restriction to a scheme etale over it is quasi-coherent. Supplies the criterion by which the six properties are transported from the cover back to the space.
\item Restriction along an etale morphism commutes with all limits and colimits. Modules on Sites, the section on localization showing that restriction has both adjoints, hence commutes with limits and colimits; Same chapter, the lemma that the pullback of a quasi-coherent module is quasi-coherent. Supplies the commutation of the restriction functor with the operations appearing in the six properties and the commutativity of the pullback diagram.
\item Commutation of the internal hom with localization. Modules on Sites, the lemma that formation of internal hom commutes with restriction. Supplies the identity used in the last of the six properties.
\end{itemize}\medskip
\ifdefined\hypertarget\hypertarget{sourcepointer-1}{}\fi
\noindent\textit{Source reference.} Original lemma 0GNB: sources/stacks-chow/tags.txt, source lines 20786--20786. Prerequisite pointer invoked at proof line(s) 5684; it names a standard fact of a sibling chapter that the counted argument uses, and every step of the counted proof built on it is retained verbatim in the delivered body.\par\medskip
\ifdefined\hypertarget\hypertarget{sourcepointer-2}{}\fi
\noindent\textit{Source reference.} Original equation 04M1: sources/stacks-spaces-properties/pinned-original-source.tex, source lines 4733--4733. Prerequisite pointer invoked at proof line(s) 5768; it names a standard fact of the same chapter that the counted argument uses, and every step of the counted proof built on it is retained verbatim in the delivered body.\par\medskip
\ifdefined\hypertarget\hypertarget{sourcepointer-3}{}\fi
\noindent\textit{Source reference.} Original equation 03LW: sources/stacks-spaces-properties/pinned-original-source.tex, source lines 4629--4629. Prerequisite pointer invoked at proof line(s) 5704; it names a standard fact of the same chapter that the counted argument uses, and every step of the counted proof built on it is retained verbatim in the delivered body.\par\medskip
\ifdefined\hypertarget\hypertarget{sourcepointer-4}{}\fi
\noindent\textit{Source reference.} Original reference 03M0: sources/stacks-spaces-properties/pinned-original-source.tex, source lines 5609--5609. Prerequisite pointer invoked at proof line(s) 5712; it names a standard fact of the same chapter that the counted argument uses, and every step of the counted proof built on it is retained verbatim in the delivered body.\par\medskip
\ifdefined\hypertarget\hypertarget{sourcepointer-5}{}\fi
\noindent\textit{Source reference.} Original reference 01BG: sources/stacks-spaces-properties/pinned-original-source.tex, source lines 5461--5461. Prerequisite pointer invoked at proof line(s) 5694; it names a standard fact of the same chapter that the counted argument uses, and every step of the counted proof built on it is retained verbatim in the delivered body.\par\medskip
\ifdefined\hypertarget\hypertarget{sourcepointer-6}{}\fi
\noindent\textit{Source reference.} Original section 00XZ: sources/stacks-spaces-properties/pinned-original-source.tex, source lines 4692--4692. Prerequisite pointer invoked at proof line(s) 5767; it names a standard fact of the same chapter that the counted argument uses, and every step of the counted proof built on it is retained verbatim in the delivered body.\par\medskip
\ifdefined\hypertarget\hypertarget{sourcepointer-7}{}\fi
\noindent\textit{Source reference.} Original section 05VR: sources/stacks-spaces-properties/pinned-original-source.tex, source lines 5784--5784. Prerequisite pointer invoked at proof line(s) 5675; it names a standard fact of the same chapter that the counted argument uses, and every step of the counted proof built on it is retained verbatim in the delivered body.\par\medskip
\ifdefined\hypertarget\hypertarget{sourcepointer-8}{}\fi
\noindent\textit{Source reference.} Original lemma 0E8H: sources/stacks-chow/tags.txt, source lines 17817--17817. Prerequisite pointer invoked at proof line(s) 5770; it names a standard fact of a sibling chapter that the counted argument uses, and every step of the counted proof built on it is retained verbatim in the delivered body.\par\medskip
\ifdefined\hypertarget\hypertarget{sourcepointer-9}{}\fi
\noindent\textit{Source reference.} Original section 03DH: sources/stacks-chow/tags.txt, source lines 4420--4420. Prerequisite pointer invoked at proof line(s) 5707; it names a standard fact of a sibling chapter that the counted argument uses, and every step of the counted proof built on it is retained verbatim in the delivered body.\par\medskip
\begin{lemma}
\label{lemma-properties-quasi-coherent}
Let $S$ be a scheme. Let $X$ be an algebraic space over $S$.
The category $\QCoh(\mathcal{O}_X)$ of quasi-coherent sheaves on $X$
has the following properties:
\begin{enumerate}
\item Any direct sum of quasi-coherent sheaves is quasi-coherent.
\item Any colimit of quasi-coherent sheaves is quasi-coherent.
\item The kernel and cokernel of a morphism of quasi-coherent sheaves
is quasi-coherent.
\item Given a short exact sequence of $\mathcal{O}_X$-modules
$0 \to \mathcal{F}_1 \to \mathcal{F}_2 \to \mathcal{F}_3 \to 0$
if two out of three are quasi-coherent so is the third.
\item Given two quasi-coherent $\mathcal{O}_X$-modules
the tensor product is quasi-coherent.
\item Given two quasi-coherent $\mathcal{O}_X$-modules
$\mathcal{F}$, $\mathcal{G}$ such that $\mathcal{F}$
is of finite presentation (see
Section \ref{section-properties-modules}),
then the internal hom
$\SheafHom_{\mathcal{O}_X}(\mathcal{F}, \mathcal{G})$
is quasi-coherent.
\end{enumerate}
\end{lemma}

\begin{proof}
If $X$ is a scheme, then this is
Descent, Lemma \ref{descent-lemma-properties-quasi-coherent}.
We will reduce the lemma to this case by \'etale localization.

\medskip\noindent
Choose a scheme $U$ and a surjective \'etale morphism $\varphi : U \to X$.
Our notation will be that $\textit{Mod}(\mathcal{O}_U) =
\textit{Mod}(U_\etale, \mathcal{O}_U)$ and
$\QCoh(\mathcal{O}_U) = \QCoh(U_\etale, \mathcal{O}_U)$; in other
words, even though $U$ is a scheme we think of quasi-coherent
modules on $U$ as modules on the small \'etale site of $U$.
By Lemma \ref{lemma-pullback-quasi-coherent} we have a commutative diagram
$$
\xymatrix{
\QCoh(\mathcal{O}_X) \ar[r]_{\varphi^*} \ar[d] &
\QCoh(\mathcal{O}_U) \ar[d] \\
\textit{Mod}(\mathcal{O}_X) \ar[r]^{\varphi^*} &
\textit{Mod}(\mathcal{O}_U)
}
$$
The bottom horizontal arrow is the restriction functor
(\ref{equation-restrict-modules})
$\mathcal{G} \mapsto \mathcal{G}|_{U_\etale}$.
This functor has both a left adjoint and a right adjoint, see
Modules on Sites, Section \ref{sites-modules-section-localize},
hence commutes with all limits and colimits.
Moreover, we know that an object of $\textit{Mod}(\mathcal{O}_X)$ is in
$\QCoh(\mathcal{O}_X)$ if and only if its restriction to $U$ is in
$\QCoh(\mathcal{O}_U)$, see
Lemma \ref{lemma-characterize-quasi-coherent}.
With these preliminaries out of the way we can start the proof.

\medskip\noindent
Proof of (1). Let $\mathcal{F}_i$, $i \in I$ be a family of quasi-coherent
$\mathcal{O}_X$-modules. By the discussion above we have
$$
\Big(\bigoplus \mathcal{F}_i\Big)|_{U_\etale} =
\bigoplus \mathcal{F}_i|_{U_\etale}
$$
Each of the modules $\mathcal{F}_i|_{U_\etale}$ is quasi-coherent.
Hence the direct sum is quasi-coherent by the case of schemes.
Hence $\bigoplus \mathcal{F}_i$ is quasi-coherent as a module
restricting to a quasi-coherent module on $U$.

\medskip\noindent
Proof of (2). Let $\mathcal{I} \to \QCoh(\mathcal{O}_X)$,
$i \mapsto \mathcal{F}_i$ be a diagram. Then
$$
(\colim \mathcal{F}_i)|_{U_\etale} = \colim \mathcal{F}_i|_{U_\etale}
$$
by the discussion above and we conclude in the same manner.

\medskip\noindent
Proof of (3). Let $a : \mathcal{F} \to \mathcal{F}'$
be an arrow of $\QCoh(\mathcal{O}_X)$. Then
we have $\Ker(a)|_{U_\etale} = \Ker(a|_{U_\etale})$
and $\Coker(a)|_{U_\etale} = \Coker(a|_{U_\etale})$
and we conclude in the same manner.

\medskip\noindent
Proof of (4). The restriction
$0 \to \mathcal{F}_1|_{U_\etale} \to \mathcal{F}_2|_{U_\etale}
\to \mathcal{F}_3|_{U_\etale} \to 0$
is short exact. Hence we have the 2-out-of-3 property
for this sequence and we conclude as before.

\medskip\noindent
Proof of (5). Let $\mathcal{F}$ and $\mathcal{G}$ be in $\QCoh(\mathcal{O}_X)$.
Then we have
$$
(\mathcal{F} \otimes_{\mathcal{O}_X} \mathcal{G})_{U_\etale} =
\mathcal{F}|_{U_\etale} \otimes_{\mathcal{O}_U} \mathcal{G}|_{U_\etale}
$$
and we conclude as before.

\medskip\noindent
Proof of (6). Let $\mathcal{F}$ and $\mathcal{G}$
be in $\QCoh(\mathcal{O}_X)$ with $\mathcal{F}$
of finite presentation. We have
$$
\SheafHom_{\mathcal{O}_X}(\mathcal{F}, \mathcal{G})|_{U_\etale} =
\SheafHom_{\mathcal{O}_U}(\mathcal{F}|_{U_\etale}, \mathcal{G}|_{U_\etale})
$$
Namely, restriction is a localization, see
Section \ref{section-localize}, especially formula
(\ref{equation-localize-at-scheme-ringed})) and formation of internal
hom commutes with localization, see
Modules on Sites, Lemma \ref{sites-modules-lemma-internal-hom-restriction}.
Thus we conclude as before.
\end{proof}

\end{document}
